\documentclass[a4paper,11pt]{article}
\usepackage{fontspec}
\setmainfont{Inter}
\setsansfont{Inter}
\renewcommand{\familydefault}{\sfdefault}
\usepackage[top=0.9cm, bottom=0.7cm, left=1.5cm, right=1.5cm]{geometry}
\usepackage{enumitem}
\usepackage{xcolor}
\usepackage{hyperref}
\definecolor{navy}{RGB}{13,44,110}
\hypersetup{
    colorlinks=true,
    urlcolor=navy,
    linkcolor=navy,
    pdftitle={Currículo — Christiano Gonçalves},
    pdfauthor={Christiano Gonçalves}
}
\linespread{1.08}

\newcommand{\secao}[1]{%
  \par
  \vspace{7pt}
  \noindent\hspace{-\fboxsep}\colorbox{black}{\makebox[\linewidth][l]{%
    \hspace{4pt}\small\bfseries\color{white}\uppercase{#1}%
  }}%
  \par\vspace{3pt}
}

\newcommand{\entrada}[4]{%
  \textbf{#1} \hfill \textit{#2}\par
  \ifx\relax#3\relax\else{\small\textit{\color{gray}#3}\par}\fi
  #4\par
  \vspace{6pt}
}

\pagestyle{empty}
\parindent=0pt
\parskip=0pt
\begin{document}

\begin{center}
{\LARGE\bfseries Christiano Gonçalves}\par\vspace{3pt}
{\normalsize Estudante de Engenharia de Software | Full Stack e Dados}\par\vspace{4pt}
{\small christiano.gonara@gmail.com $\cdot$ (31) 99890-9804 $\cdot$ Belo Horizonte, MG\par}
{\small
\href{https://linkedin.com/in/christiano-gonara}{linkedin.com/in/christiano-gonara}
$\cdot$
\href{https://github.com/christiano-gonara}{github.com/christiano-gonara}\par}
\end{center}

\secao{Objetivo}
Graduando em Engenharia de Software na PUC Minas, com experiência em automação e dados com Python e em aplicações web com React. Busco estágio em desenvolvimento web ou dados.\par

\secao{Habilidades}
\textbf{Linguagens:} Python, JavaScript, Java, C, HTML, CSS\par
\textbf{Backend:} Flask, FastAPI, Spring Boot, REST API, Scrapy, pytest\par
\textbf{Dados e Banco:} PostgreSQL, SQLite, DuckDB, Parquet\par
\textbf{Frontend:} React, Next.js, Vite, Tailwind CSS, Bootstrap, Chart.js\par
\textbf{Ferramentas e IA:} Git, GitHub, Docker, Railway, Render, Claude Code, Codex\par
\textbf{Idiomas:} Inglês técnico (leitura de documentação)\par

\secao{Projetos}
\entrada{Desenvolvedor Backend e Dados — CNPJ Alimentação}{2026}
{Python · DuckDB · Parquet · FastAPI · React}
{Pipeline com dados abertos da Receita Federal sobre alimentação fora do lar. Processei 72,8 milhões de registros com DuckDB e construí um dashboard em React que resume o setor.}

\entrada{Desenvolvedor Backend — Padrinho Track}{Mar/2026 a Jul/2026}
{Python · Flask · PostgreSQL · Google Sheets API · pytest}
{Sistema em produção que gerencia o Programa de Apadrinhamento da PUC Minas (50 padrinhos e 127 calouros). Gera certificados e relatórios automaticamente e tem 62 testes automatizados.}

\entrada{Hackathon Localiza — 4º lugar entre 40 equipes}{2026}
{JavaScript · HTML · CSS · API externa · GitHub Pages}
{Simulador que mostra se um carro elétrico por assinatura cabe na rotina do cliente, comparando os custos com a gasolina por estado via API.}

\entrada{Desenvolvedor Full Stack — MinhaReceita}{2026}
{React · Next.js · MongoDB · Docker}
{Receituário digital para o SUS, feito em equipe de quatro: o médico prescreve, o sistema gera um código e a farmácia valida. Desenvolvi a landing page e as telas do médico em React.}

\secao{Experiência}
\entrada{Coordenador de Mentoria Acadêmica — PUC Minas}{Fev/2026 a Atual}{}
{Coordeno cerca de 50 padrinhos voluntários no acompanhamento de 127 calouros. Planejo o cronograma e conduzo as reuniões do programa.}

\secao{Formação}
\textbf{Graduação em Engenharia de Software} \hfill \textit{Jan/2025 a Dez/2028}\par
PUC Minas\par\vspace{3pt}
\textbf{Python e Análise de Dados (20h)} \hfill \textit{Mar/2026}\par
Universidade Federal de Minas Gerais (UFMG)\par\vspace{3pt}
\textbf{Claude Code in Action} \hfill \textit{Jun/2026}\par
Anthropic\par\vspace{3pt}
\textbf{Comunicação e Oratória} \hfill \textit{Mar/2026 a Dez/2026}\par
Escola de Teatro PUC Minas\par

\end{document}
